% Table 3 — HG-Rec curvature ablation vs other generative baselines
% Source: Task #246 v3 ranking + Task #88/#89 per-layer curvature analysis
\begin{table}[t]
\centering
\caption{HG-Rec curvature ablation vs other generative baselines on Musical_Instruments.
HG-Rec c555 marginally outperforms phonism (within seed noise). All HG-Rec variants
significantly outperform other generative methods (LETTER / TIGER / FDSA) by +76\% on average,
confirming the RQ-VAE + T5 generation pipeline as the dominant performance driver.
Geometric prior is marginal (5.3\% span across 6 curvature grids, Task~\#88).}
\label{tab:hgrec_vs_generative}
\small
\begin{tabular}{l r r l}
\toprule
Method & R@10 & vs phonism & Notes \\
\midrule
\textbf{RQ-VAE + T5} & & & \\
\quad phonism (vanilla + Sinkhorn) & \textbf{0.1058} & -- & best overall \\
\quad HG-Rec c555 ($\kappa$=0.5) & 0.1051 & -0.7\% & marginally below phonism \\
\quad HG-Rec c111 (paper default) & 0.1020 & -3.6\% & HG-Rec baseline \\
\midrule
\textbf{Other generative} & & & \\
\quad DuoRec & 0.0672 & -36.4\% & sequential encoder-decoder \\
\quad FDSA & 0.0594 & -43.8\% & feature-disentangled self-attn \\
\quad TIGER (T5-small) & 0.0591 & -44.1\% & quantizer-free codebook \\
\quad LETTER (paper-aligned) & 0.0509 & -51.9\% & after Task~\#150 fix \\
\midrule
\textbf{Average gain} & & & \\
\quad RQ-VAE + T5 (3 variants) & 0.1043 & -- & -- \\
\quad Other generative (4 methods) & 0.0592 & -- & -- \\
\quad \textbf{RQ-VAE + T5 advantage} & \textbf{+76\%} & -- & $0.1043/0.0592 - 1$ \\
\bottomrule
\end{tabular}
\end{table}